\clearpage\corrigeDevoir{3}
\begin{blocLivre}{dscor-3-1}\begin{corrigeBac}{6}{Une réserve de glace}
\demarche{Convertir la masse et les kilojoules avant d’utiliser le débit d’énergie.}
\begin{questions}
\item $\ce{H2O(s) -> H2O(l)}$ ; fusion endothermique.\pts{2 pt}
\item $E=0,150\times334=50,1$ kJ.\pts{2 pt}
\item $t=50100/50=1002$ s, environ 16,7 min. Une fois la glace disparue, l’eau liquide peut se réchauffer.\pts{2 pt}
\end{questions}
\end{corrigeBac}\end{blocLivre}
\begin{blocLivre}{dscor-3-2}\begin{corrigeBac}{7}{Détartrage contrôlé}
\demarche{Comparer les quantités divisées par les coefficients, puis reconstruire le bilan final.}
\begin{questions}
\item L’ion chlorure est spectateur.\pts{1 pt}
\item $0,040/1>0,060/2=0,030$ : les ions hydrogène limitent.\pts{2 pt}
\item 0,030 mol de calcaire réagissent ; 0,010 mol restent ; 0,030 mol de \ce{CO2} se forment.\pts{2 pt}
\item Le gaz trouble l’eau de chaux. Ne pas fermer hermétiquement le système qui dégage du gaz.\pts{2 pt}
\end{questions}
\end{corrigeBac}\end{blocLivre}
\begin{blocLivre}{dscor-3-3}\begin{corrigeBac}{7}{Trois transformations}
\demarche{Analyser séparément le contrôle chimique et le bilan des noyaux.}
\schema[corrige]{ccm-synthese}
\begin{questions}
\item La synthèse redistribue les atomes sans changer les noyaux ; la réaction nucléaire modifie les noyaux.\pts{1 pt}
\item Une tache de S correspond à N : présence compatible avec le produit recherché. Une seconde correspond à R : S n’est pas pur dans cette analyse.\pts{2 pt}
\item Il condense les vapeurs qui retombent dans le ballon, limitant les pertes pendant le chauffage.\pts{2 pt}
\item $A:2+3=4+1$ et $Z:1+1=2+0$ ; c’est une fusion nucléaire.\pts{2 pt}
\end{questions}
\end{corrigeBac}\end{blocLivre}

